\section{Décomposition en \(5\) modes de polarisation}

La dimension cinq acquiert un contenu tensoriel lorsque l’on choisit une direction
unitaire de propagation \(n\) et un repère orthonormé orienté
\((e_1,e_2,n)\). Nous définissons les cinq tenseurs réels
\begin{align}
E_{+}&=\frac1{\sqrt2}(e_1e_1^T-e_2e_2^T),
&E_{\times}&=\frac1{\sqrt2}(e_1e_2^T+e_2e_1^T),
\label{rm:eq:gravity-helicity2-real}\\
E_{1}&=\frac1{\sqrt2}(e_1n^T+ne_1^T),
&E_{2}&=\frac1{\sqrt2}(e_2n^T+ne_2^T),
\label{rm:eq:gravity-helicity1-real}\\
E_{0}&=\frac1{\sqrt6}(2nn^T-e_1e_1^T-e_2e_2^T).
\label{rm:eq:gravity-helicity0-real}
\end{align}
Ils sont tous symétriques et sans trace. Les deux premiers sont transversaux ; les
deux suivants mêlent le plan transversal à la direction de
propagation, et le dernier est le mode scalaire longitudinal sans trace.

\begin{theorem}[Décomposition hélicoïdale complète du porteur]
\label{rm:thm:gravity-five-helicities}
Les tenseurs de \cref{rm:eq:gravity-helicity2-real,rm:eq:gravity-helicity1-real,rm:eq:gravity-helicity0-real} forment une base orthonormale de \(\operatorname{STF}_3\). Dans la complexification, les combinaisons
\begin{equation}
\label{rm:eq:gravity-complex-helicities}
E_{\pm2}=\frac1{\sqrt2}(E_+\mp\ii E_\times),
\qquad
E_{\pm1}=\frac1{\sqrt2}(E_1\mp\ii E_2),
\qquad E_0=E_0
\end{equation}
sont des vecteurs de poids \(\pm2,\pm1,0\), respectivement, sous les rotations
autour de \(n\). De plus,
\begin{equation}
\label{rm:eq:gravity-tt-on-five}
\Lambda(n)E_{\pm2}=E_{\pm2},
\qquad
\Lambda(n)E_{\pm1}=0,
\qquad
\Lambda(n)E_0=0.
\end{equation}
\end{theorem}

\begin{proof}
Dans le repère \((e_1,e_2,n)\), chaque tenseur a une norme de Frobenius égale à un.
Leurs motifs d’entrées non diagonales sont disjoints, et les trois diagonales
de \(E_+,E_0\) ont un produit intérieur nul ; les cinq tenseurs
sont donc orthonormés. Comme \(\dim\operatorname{STF}_3=6-1=5\), ils forment une base.

Une rotation d’angle \(\psi\) autour de \(n\) envoie
\((e_1,e_2)\) sur
\((\cos\psi\,e_1+\sin\psi\,e_2,
-\sin\psi\,e_1+\cos\psi\,e_2)\). En substituant dans les formules, on obtient
une rotation d’angle \(2\psi\) sur
\(\operatorname{span}\{E_+,E_\times\}\), une d’angle \(\psi\) sur
\(\operatorname{span}\{E_1,E_2\}\), et \(E_0\) fixe. D’où les poids de
\cref{rm:eq:gravity-complex-helicities}. Enfin,
\(\Pi E_{+}\Pi=E_+\) et \(\Pi E_\times\Pi=E_\times\), tandis que
\(\Pi E_1\Pi=\Pi E_2\Pi=0\). Pour \(E_0\),
\(\Pi E_0\Pi=-(e_1e_1^T+e_2e_2^T)/\sqrt6\), qui est purement de trace dans
le plan, et la soustraction de \cref{rm:eq:gravity-lambda-tt} l’annule.
\end{proof}

Le théorème précise le sens de ``cinq polarisations'' sans leur attribuer
à toutes le même statut dynamique. Avant d’imposer les équations de champ,
le porteur de spin deux a cinq coordonnées irréductibles. Dans une
réalisation massive de Fierz--Pauli, les contraintes laissent précisément
les cinq poids de \cref{rm:eq:gravity-complex-helicities} ; dans une réalisation
sans masse avec invariance de jauge, les secteurs \(\pm1\) et \(0\) sont éliminés et
il reste \(\pm2\). HMT fournit ici le porteur, la base et le projecteur ;
le choix de la dynamique massive ou sans masse relève de l’interface
physique. Cinq ne signifie donc pas ``cinq ondes observées'', mais cinq
degrés tensoriels disponibles avant la réduction.

L’orientation reste également visible : remplacer
\((e_1,e_2,n)\) par \((e_1,-e_2,-n)\) conserve \(E_+,E_0\) et échange
les phases des paires hélicoïdales. Cette information ne se trouve pas dans le
scalaire \(G_a\) ; elle réside dans le porteur sur lequel agit
\(\mathbb G_{5,a}\).

\section{Chaîne représentationnelle \(12\to5\to5\to2\) : support \(V_5\), tenseur sans trace, polarisations de Fierz–Pauli et hélicités transverses}

La dimension cinq apparaît dans des objets distincts, et seul
\cref{rm:eq:gravity-12-5-5-2} autorise le transport entre eux :

\begin{enumerate}
\item \(V_5\) est le facteur irréductible de la représentation sur les
douze sommets. C’est un résultat fini de théorie des représentations.
\item \(\operatorname{STF}_3\) est le porteur réel de dimension cinq des
poids \(m=-2,-1,0,1,2\) de spin deux.
\item Une théorie massive de Fierz--Pauli peut réaliser ces cinq poids
comme cinq états physiques, après fixation de l’action et des contraintes.
\item La relativité générale sans masse conserve deux hélicités radiatives,
\(\pm2\), représentées par \(E_+\) et \(E_\times\).
\item Les classes \(E(2)\) d’ondes métriques, comme \(III_5\), classent
les amplitudes possibles sous d’autres hypothèses ; elles ne sont pas automatiquement
\(V_5\), Fierz--Pauli ou la relativité générale.
\end{enumerate}

En particulier, la formulation correcte est : \emph{HMT construit un support tensoriel à cinq composantes avant les contraintes ; la réduction TT de la réalisation sans masse est de rang deux}. Il n’est pas affirmé que la relativité générale possède cinq polarisations propagatives.
